\begin{abstract}
La convergencia de la computación cuántica y las redes Narrowband IoT (NB-IoT) presenta desafíos significativos para la ciberseguridad en entornos de exploración profunda, donde los dispositivos de bajo consumo transmiten datos sensibles en entornos remotos y hostiles. Este artículo analiza las vulnerabilidades de los esquemas criptográficos clásicos frente a algoritmos cuánticos como Shor y Grover, y propone un marco de seguridad cuántica basado en criptografía post-cuántica (PQC) y distribución de claves cuánticas (QKD) adaptada a las restricciones de NB-IoT. Se evalúa el estado del arte en implementaciones PQC para dispositivos restringidos, se presenta una arquitectura híbrida PQC-aware para onboarding en plano de control y protección ligera en plano de datos, y se analizan métricas de rendimiento como latencia, consumo energético y overhead de comunicación. Los resultados demuestran viabilidad práctica con optimizaciones selectivas, manteniendo compatibilidad con estándares 3GPP mientras se logra resistencia cuántica. Se discuten implicaciones para aplicaciones críticas como monitoreo ambiental y exploración industrial.
\end{abstract}